\documentclass[11pt]{article}
\usepackage{amsmath,amssymb,amsthm}
\newtheorem{theorem}{Theorem}
\newtheorem{lemma}[theorem]{Lemma}
\title{Moment Fibers and External Stieltjes Readouts}
\author{}
\date{November 23, 2026}
\begin{document}
\maketitle

\paragraph{Status and scope.}
These are complete written arguments, with exact small checks in the
companion Python file. They have not been formalized in Lean. They audit
a possible mechanism for the main asymptotic target: preserving many
moments of a positive atomic measure does not force its remaining
Stieltjes differences to have one sign at all admissible query poles.
The independent readout theorem below has only linear efficiency and
does not prove either requested global bound.

\begin{lemma}[Reciprocal powers]
For distinct real $t_1,\ldots,t_q$, a nonzero linear combination of
$(t_i-v)^{-k}$, where $k$ is a positive integer, has at most $q-1$
distinct zeros on $v<\min_i t_i$.
\end{lemma}
\begin{proof}
Replace $v$ by $-v$. It suffices to consider
$\sum_j a_j(u-\lambda_j)^{-k}$ on $u>\max_j\lambda_j$,
with $\lambda_1<\cdots<\lambda_q$.
Induct on $q$, simultaneously for all $k$. If the sum had $q$ zeros,
multiply by $(u-\lambda_q)^k$ and differentiate. Rolle's theorem
would give $q-1$ zeros of
\[
 k(u-\lambda_q)^{k-1}
 \sum_{j<q} a_j(\lambda_q-\lambda_j)(u-\lambda_j)^{-k-1}.
\]
The induction hypothesis allows at most $q-2$ unless all $a_j$ for
$j<q$ vanish. In the latter case the original nonzero sum has no
zero. The case $q=1$ is immediate.
\end{proof}

\begin{theorem}[Independent readouts while preserving moments]
Let $K,q\geq1$ and $H=K+q$. Fix distinct positive numbers
$\lambda_1,\ldots,\lambda_H$ and distinct query poles
$t_1,\ldots,t_q>\max_h\lambda_h$. For every sufficiently small
vector $b\in\mathbb R^q$, there exist distinct positive
$v_1,\ldots,v_H$, arbitrarily close to the corresponding
$\lambda_h$ as $b\to0$, such that
\[
 \sum_h v_h^j=\sum_h\lambda_h^j\quad(1\leq j\leq K),
 \qquad
 \sum_h\frac1{t_i-v_h}-\sum_h\frac1{t_i-\lambda_h}=b_i
 \quad(1\leq i\leq q).
 \tag{1}
\]
All $v_h$ remain below every query pole. In particular all $2^q$
strict sign patterns of the readout differences are possible while
these $K$ moments are fixed.
\end{theorem}
\begin{proof}
Apply the inverse function theorem to the map from the $H$ atom
positions to the $K$ power sums and $q$ Stieltjes values. Its
Jacobian rows at $\lambda$ are
\[
 (j\lambda_h^{j-1})_h\quad(1\leq j\leq K),
 \qquad ((t_i-\lambda_h)^{-2})_h\quad(1\leq i\leq q).
 \tag{2}
\]
A row dependence would give a polynomial $p$ of degree at most
$K-1$ and coefficients $c_i$ such that
\[
       p(v)+\sum_i c_i(t_i-v)^{-2}
 \tag{3}
\]
vanishes at the $H$ distinct atom positions. Differentiating $K$
times and applying Rolle's theorem repeatedly gives at least
$H-K=q$ distinct zeros of
\[
               (K+1)!\sum_i c_i(t_i-v)^{-K-2}.
\]
The lemma forces every $c_i$ to vanish. Then $p$ has $H>K-1$
zeros, so it too vanishes. Thus (2) is invertible. Positivity,
distinctness, and the strict gaps from the query poles are open
conditions, so the inverse function theorem can be applied inside
their common domain. Set the moment components of its target equal
to their original values and vary only the Stieltjes components.
\end{proof}

\paragraph{Transfer to strictly admissible cube heads.}
Take $q=2^d$ for $d\geq1$ and choose the query poles to be the
distinct values of
\[
       t(y)=C-\sum_{j=0}^{d-1}2^j y_j,
       \qquad y\in\{0,1\}^d,
 \tag{4}
\]
where $C-(q-1)>\max_h\lambda_h$.
For a small target vector in (1), fix
$a=\min_h\lambda_h/4$ and cyclically define
\[
 w_1=a,\qquad w_{h+1}=w_h+v_h-\lambda_h,
 \qquad u_h=\lambda_h-w_h.
 \tag{5}
\]
The first-moment equality ensures $w_{H+1}=w_1$. By taking the
target vector sufficiently small, all $w_h$ and $u_h$ are positive.
Moreover
\[
       u_h+w_h=\lambda_h,
       \qquad u_h+w_{h+1}=v_h.
 \tag{6}
\]
The following score uses $2H$ actual affine-ratio heads on the
unchanged $(d+2)$-dimensional cube:
\[
 G(x_1,x_2,y)=
 \sum_{h=1}^H\frac1{t(y)-u_hx_1-w_hx_2}
 -\sum_{h=1}^H\frac1{t(y)-u_hx_1-w_{h+1}x_2}.
 \tag{7}
\]
Every denominator slope is strictly negative. Its minimum on the
whole cube is either $\min_y t(y)-\lambda_h>0$ or
$\min_y t(y)-v_h>0$, so the original atom admissibility applies
directly. When $x_1=0$, cyclic permutation of the $w_h$ cancels the
two sums; when $x_2=0$, the cancellation is termwise. At
$x_1=x_2=1$, equations (1) and (6) give
\[
                 G(1,1,y)=-b_y.
 \tag{8}
\]
Thus arbitrary strict signs on this codimension-two face are
possible while the score vanishes on its two coordinate axes and
the first $K$ spectral moments agree. The budget is $2(K+2^d)$
heads, not a superconstant improvement over the number of face
points. No asymptotic saving is being asserted.

\paragraph{An exact two-query example.}
The spectra $A=(3,4,9)$ and $B=(1,7,8)$ have equal mass and equal
mean. Their Stieltjes difference is
\[
 R(t)=\sum_{a\in A}\frac1{t-a}-\sum_{b\in B}\frac1{t-b}
 =\frac{-4(2t^3-55t^2+416t-867)}
 {(t-9)(t-8)(t-7)(t-4)(t-3)(t-1)}.
 \tag{9}
\]
Although both query poles lie strictly above all atoms,
\[
                 R(10)=\frac{23}{63}>0,
 \qquad R(20)=-\frac{1453}{2217072}<0.
 \tag{10}
\]
For a direct cube realization, take
$u=(1/2,7/2,11/2)$ and $w=(5/2,1/2,7/2)$.
Then $u_h+w_h$ gives $A$, while $u_h+w_{h+1}$ gives $B$,
with cyclic indexing. Set $t(y)=20-10y$ in (7). The score has
six heads on three actual Boolean coordinates, positive cube
denominators, strictly negative slopes, exact axis cancellation,
and opposite strict signs at the two mixed-face query vertices.

\paragraph{Crossing after almost all moments are fixed.}
For every $K\geq1$, set $H=K+2$ and
\[
 P(t)=\prod_{j=1}^H(t-j),\qquad a=H+1,
 \qquad Q(t)=P(t)+\varepsilon(t-a).
 \tag{11}
\]
For sufficiently small $\varepsilon>0$, the polynomial $Q$ has
one distinct positive real root in each interval
$(j-1/4,j+1/4)$. This follows by preserving the opposite signs
of $P$ at the two endpoints of each interval. Newton's identities
show that the root multisets of $P$ and $Q$ have identical power
sums through degree $H-2=K$, since their coefficients of degrees
$H-1$ through $2$ coincide.
The difference of their Stieltjes sums is
\[
 \frac{P'(t)}{P(t)}-\frac{Q'(t)}{Q(t)}
 =\frac{\varepsilon((t-a)P'(t)-P(t))}{P(t)Q(t)}.
 \tag{12}
\]
Both $a$ and $2a$ exceed all roots. At $a$ this expression is
$-\varepsilon/P(a)<0$. At $2a$ its numerator is positive because
\[
 (2a-a)\frac{P'(2a)}{P(2a)}
 =a\sum_{j=1}^H\frac1{2a-j}>H/2>1.
\]
Thus equality of the first $K$ moments does not force a common
external-pole sign even with only $K+2$ atoms on each side.
The independent-readout theorem proves more, but (11)--(12)
give a particularly explicit crossing.

\paragraph{What remains missing.}
For the global target, many permutation configurations sharing
low moments might supply many branches of the same continuous
interpolation problem. Neither the number of configurations nor
the independent readouts in (1) prove that those branches cover
all required sign orthants. The present theorem provides exactly
$H-K$ independently controlled values using $H$ moving atom
positions. A superconstant improvement would require an additional
argument extracting more independent information from the branches
while preserving the original cube admissibility and total head
budget. That argument remains open.

\end{document}
